\documentclass[11pt]{article}
\usepackage{fontspec}
\usepackage[french]{babel}
\usepackage{amsmath}
\usepackage{amssymb}
\usepackage{amsthm}
\usepackage{mathrsfs}
\usepackage{tikz-cd}
\usepackage{enumitem}

\newcommand{\printedpage}[1]{\marginpar{\footnotesize p.~#1}}
\newcommand{\msp}[1]{\marginpar{\footnotesize [#1]}}
\newcommand{\Ob}{\operatorname{Ob}}
\newcommand{\Fl}{\operatorname{Fl}}
\newcommand{\Cat}{\mathrm{Cat}}
\newcommand{\Fibdisc}{\underline{\mathrm{Fib\ disc}}}
\newcommand{\cercle}[1]{\tikz[baseline=(c.base)]\node[draw,circle,inner sep=1pt](c){#1};}

\begin{document}

\printedpage{108}
% head: A. Grothendieck : Dérivateurs, Ch. XII, Édition des Universités de Montpellier II et Paris VII
(\footnote[107]{attention, ce n'est pas un carré commutatif. Dire plutôt que $(\mathcal{X}/x_0)_u \simeq \int \begin{tikzcd}[ampersand replacement=\&, row sep=0.4em, column sep=0.4em, baseline=(\tikzcdmatrixname-2-1.base)] \& {} \\ {} \ar[ur]\ar[dr] \& \\ \& {} \end{tikzcd}$ est lisse sur $\Phi$ à fibres $W$-asphériques (si $f_0$, $f_1$, $f_2$ sont lisses), donc $W$-asphérique puisque $\Phi$ l'est.}), il faut exprimer que $(\mathcal{X}/x_0)_u$ est $W$-asphérique, sachant que $(X_1/x_1)_u$ et $(X_2/x_2)_u$ le sont. Donc cela équivaut à % ?? uncertain: small diagram in footnote 107 -- two arrows (up-right and down-right) printed with no visible objects or labels
\[
i_1(x_0,u) \in W \qquad (\text{ou } i_2(x_0,u) \in W),
\]
donc aussi (moyennant W(7 bis)) à
\msp{page 157}
\[
\alpha_1(x_0,u) \in W, \qquad \text{i.e. } (X_0/x_0)_u \ W\text{-asphérique.}
\]
Dire que c'est vrai pour tout $u$ dans $Y/y$, signifie que $f_0 : X_0 \longrightarrow Y$ est lisse en $u$ \texttt{[plutôt en $x_0$]}.

Donc on a prouvé a), sous la forme plus précise :

\textbf{Corollaire 6.} \textit{Soit $x \in \Ob \mathcal{X}$ (où $\mathcal{X} = \int \begin{tikzcd}[ampersand replacement=\&, row sep=tiny, column sep=small, baseline=(\tikzcdmatrixname-2-1.base)] \& X_1 \\ X_0 \ar[ur]\ar[dr] \& \\ \& X_2 \end{tikzcd}$.) On veut exprimer la $W$-lissité de}
\[
\varphi_Q : \mathcal{X} \longrightarrow Y
\]
\textit{en $x$.}

\begin{enumerate}[label={}, leftmargin=2.5em]
\item[\textit{a)}] \textit{Si $x \in \Ob X_i$ ($i \in \{1,2\}$), $\varphi_Q$ est $W$-lisse en $x$ si et seulement si $f_i : X_i \longrightarrow Y$ l'est.}
\item[\textit{b)}] \textit{Supposons} \textup{W(7 bis)} (\footnote[108]{Inutile pour la partie directe, en supposant \texttt{[phrase incomplète]}.})\textit{. Soit $x_0 \in \Ob X_0$, et supposons $f_1$, $f_2$ lisses en $x_1$, $x_2$. Alors $\varphi_Q$ est $W$-lisse en $x_0$ si et seulement si $f_0 : X_0 \longrightarrow Y$ l'est.}
\end{enumerate}

\textit{Donc} (moyennant W(7 bis)) \textit{dans les trois cas $x \in \Ob X_i$ ($i \in \{0,1,2\}$) on trouve que $\varphi_Q$ est $W$-lisse en $x$ si et seulement si $f_i : X_i \longrightarrow Y$ l'est} \texttt{[pas vraiment pour $i=0$]}.

\cercle{b} Quant à b) du théorème 4, c'est le lemme 2 ci-dessus. Donc tout est prouvé!

\msp{page 158}

\subsection*{3.5\quad Applications aux sommes amalgamées dans une catégorie $X^{\wedge}$ (avec $X \in \Ob \Cat$).}

D'abord une caractérisation des sommes amalgamées impliquant une immersion ouverte :

\textbf{Proposition 15.} \textit{Soit}
% foot: Version 7 décembre 2002
\printedpage{109}
% head: A. Grothendieck : Dérivateurs, Ch. XII, Édition des Universités de Montpellier II et Paris VII
\[
Q : \qquad
\begin{tikzcd}
U \ar[r, "i"] \ar[d, "p"'] & X \ar[d, "q"] \\
U' \ar[r, "i'"] & X'
\end{tikzcd}
\]
\textit{un carré commutatif dans} $\Cat$\textit{. Pour que ce carré soit cocartésien dans} $\Cat$\textit{, avec $i$ (donc $i'$) une immersion ouverte, il faut et il suffit que les conditions suivantes soient satisfaites :}

\begin{enumerate}[label={}, leftmargin=2.5em]
\item[\textit{a)}] \textit{$Q$ est cartésien.}
\item[\textit{b)}] \textit{$i'$ est une immersion ouverte (donc $i$ aussi).}
\item[\textit{c)}] \textit{Si $T' = X' \backslash U'$, et $T = X \times_{X'} T' = X \backslash U$, $T \longrightarrow T'$ est un isomorphisme.}
\item[\textit{d)}] \textit{Pour tout $t \in T$, si $t'$ est son image dans $T'$, le foncteur induit par $p$}
\[
p/t : \underbrace{U/t}_{= V_t} \longrightarrow \underbrace{U'/t'}_{= V'_{t'}}
\]
\textit{satisfait à $(p/t)_!(e_{U/t}) \simeq e_{U'/t'}$, ou ce qui revient au même, les `fibres colocales' de $p/t$, $v' \backslash V_t$, avec $v'$ dans $V'_{t'}$, sont} $0$\textit{-connexes.}
\end{enumerate}

\msp{page 159}
Démonstration immédiate - ça devrait plutôt remonter au §1.

\textbf{Corollaire 1.} \textit{Soit $Q$ un carré} cartésien (cf. page 158), \textit{avec $i'$ une} immersion ouverte \textit{(donc $i$ aussi), et $q$} $W$-lisse \textit{(donc $p$ aussi). On suppose que $W$ satisfait l'axiome} W(5) \textit{de connexité} (\footnote[109]{Donc il suffit $q$ $W_0$-lisse, où $W_0 = \{f \in \Fl \Cat \mid \pi_0(f) \text{ bijectif}\}$.})\textit{. Soit $T' = X' \backslash U'$, $T = X \backslash U$, supposons que $T' \longrightarrow T$ soit un isomorphisme. Alors $Q$ est cocartésien.}

Il suffit de prouver la condition d) de prop. 15. Considérons le carré
\[
\begin{tikzcd}
U/t \ar[r, hook, "i/t"] \ar[d, "p/t"'] & X/t \ar[d, "q/t"] \\
U'/t' \ar[r, hook, "i'/t'"] & X'/t',
\end{tikzcd}
\tag*{$(*)$}
\]
je dis qu'il est \textit{cartésien.} Cela résulte de la seule hypothèse que $Q$ soit cartésien, sans hypothèse de lissité sur $q$ ni que $i'$ soit une immersion ouverte, et se voit p. ex. en regardant le diagramme cubique :
% foot: Version 7 décembre 2002
\printedpage{110}
% head: A. Grothendieck : Dérivateurs, Ch. XII, Édition des Universités de Montpellier II et Paris VII
\[
\begin{tikzcd}[row sep=small, column sep=small]
& U/t \ar[rr] \ar[dl] \ar[dd] & & X/t \ar[dl] \ar[dd] \\
U'/t' \ar[rr, crossing over] \ar[dd] & & X'/t' & \\
& U \ar[rr] \ar[dl] & & X \ar[dl] \\
U' \ar[rr] & & X' \ar[from=uu, crossing over] &
\end{tikzcd}
\]
(\footnote[110]{NB Ici $t$ est un élément quelconque de $\Ob X$, $t'$ son image dans $X'$.}). Dans ce diagramme, la face antérieure et postérieure du cube est cartésienne, par définition de $U/t$ et de $U'/t'$
\msp{page 160}
respectivement, et la base l'est par hypothèse. Il en résulte que la face supérieure du cube est également un carré cartésien.

Ceci vu, commme $q$ est lisse, il en est de même de $q/t$ dans le carré cartésien $(*)$, donc aussi de $p/t$ par changement de base. D'autre part, la lissité de $q$ implique que $q/t$ est à fibres $W$-asphériques, donc par changement de base il en est de même de $p/t$ (qui est donc une $W$-équivalence universelle, i.e. une $W$-fibration triviale). De plus, à cause de la lissité de $p/t$, pour $v'$ dans $U'/t'$, l'inclusion % [sic] commme -> comme, repeated letter
\[
(U/t)_{v'} \longrightarrow v' \backslash (U/t)
\]
est dans $W$, donc (par W(5)) définit une bijection sur les $\pi_0$. Comme $(U/t)_{v'}$ est $W$-asphérique par ce qui précède, il s'ensuit (par W(5)) que son $\pi_0$ est de cardinal 1, donc il en est de même de $v' \backslash (U/t)$, i.e. celui-ci est 0-connexe. Or c'est là la condition d) de la proposition 15, qu'il fallait établir.

\msp{page 161}
\textbf{Corollaire 2.} \textit{Soit $Q$ un carré cartésien} (p. 158), \textit{avec $i'$ (donc aussi $i$) une immersion ouverte, et $q$ (donc aussi $p$) un} isomorphisme local. \textit{Soit $T = X \backslash U$, $T' = X' \backslash U'$ (donc $T \simeq X \times_{X'} T'$), et supposons $T' \longrightarrow T$ isomorphisme. Alors $Q$ est cocartésien.}

C'est un cas particulier du corollaire 1, mais sans référence explicite à un localiseur $W$. On peut pendre $W = W_{\infty}$, ou $W = W_0$, pour se ramener au corollaire 1. Mais il est plus raisonnable de dire que dans le carré cartésien $(*)$, l'hypothèse que $q$ soit un isomorphisme local implique que $q/t$ est un isomorphisme, donc $p/t$ l'est aussi. Il en résulte bien que les $v' \backslash (U/t)$ sont 0-connexes, cette catégorie étant isomorphe à $v' \backslash (U'/t')$, donc a un objet initial. % [sic] pendre -> prendre, missing letter

Soit maintenant $X$ un objet de Cat, et considérons un diagramme cocartésien dans $X^{\wedge}$ :
\[
\begin{tikzcd}
G \ar[r, "i"] \ar[d, "p"'] & F \ar[d, "q"] \\
G' \ar[r, "i'"] & F',
\end{tikzcd}
\tag*{$(*)$}
\]
% foot: Version 7 décembre 2002
\printedpage{111}
% head: A. Grothendieck : Dérivateurs, Ch. XII, Édition des Universités de Montpellier II et Paris VII
\msp{page 162}
où $i$ est un monomorphisme. On sait que le caractère cocartésien du carré se vérifie `argument par argument' par le caractère cocartésien de
\[
\begin{tikzcd}
G(x) \ar[r, "i(x)"] \ar[d, "p(x)"'] & F(x) \ar[d, "q(x)"] \\
G'(x) \ar[r, "i'(x)"] & F'(x),
\end{tikzcd}
\tag*{$(*)$}
\]
dans Ens, $\forall\, x$ dans $X$. Or un tel carré cocartésien dans Ens, avec $i(x)$ \texttt{[un]} monomorphisme, est aussi cartésien. Cela implique que le carré $Q$ dans $X^{\wedge}$ est également cartésien.

Considérons, d'autre part, le foncteur canonique
\[
\varphi_X : F \longmapsto X/F : X^{\wedge} \longrightarrow \mathrm{Cat}.
\]
Il vaut mieux, d'ailleurs, le regarder comme un foncteur
\[
\varphi_X : X^{\wedge} \longrightarrow \mathrm{Cat}/X,
\]
et je rappelle, pour mémoire, que ce foncteur est pleinement fidèle, et que son image essentielle est formée des catégories $X'$ sur $X$ telles que
\[
X' \longrightarrow X
\]
\msp{page 163}
soit un isomorphisme local, ou ce qui revient au même, que $X'$ soit \textit{fibrée} sur $X$ à fibres discrètes :
\[
\varphi_X : X^{\wedge} \xrightarrow{\ \sim\ } \underbrace{\Fibdisc\, X}_{\substack{\text{catégories fibrées}\\ \text{à fibres discrètes}}} \ .
\]
Notons que toute flèche $X' \longrightarrow X''$ dans $\Fibdisc\, X$ est un isomorphisme local (tautologique sur $\begin{tikzcd}[ampersand replacement=\&, row sep=small, column sep=tiny, baseline=(\tikzcdmatrixname-1-1.base)] X'/x' \ar[rr] \ar[dr, "\text{iso.}"'] \& \& X''/x'' \ar[dl, "\text{iso.}"] \\ \& X/x \& \end{tikzcd}$), ce qui montre que si on a un diagramme
\[
\begin{tikzcd}[row sep=small, column sep=small]
X_1 \ar[dr] & & X_2 \ar[dl] \\
& X_0 &
\end{tikzcd}
\]
dans $\Fibdisc\, X$, le produit fibré $X_3 = X_1 \times_{X_0} X_2$ dans Cat est tel que $X_3 \longrightarrow X_2$ est un isomorphisme local, donc aussi $X_3 \longrightarrow X$. C'est-à-dire $\Fibdisc(X)$ est stable par produits fibrés. Ceci implique que les produits fibrés existent dans $\Fibdisc\, X$, et que ce sont les produits fibrés dans Cat, ou dans Cat$/X$.
% foot: Version 7 décembre 2002

\printedpage{112}
% head: A. Grothendieck : Dérivateurs, Ch. XII, Édition des Universités de Montpellier II et Paris VII
\msp{page 164}
De même pour les produits finis dans $\Fibdisc\, X$, ce sont les produits finis dans Cat$/X$. Bref

\textbf{Lemme 1.} \textit{La sous-catégorie pleine} $\Fibdisc\, X$ \textit{de} $\Cat/X$ \textit{est stable par} $\varprojlim$ \textit{finies, donc les} $\varprojlim$ \textit{existent dans} $\Fibdisc\, X$ \textit{et le foncteur d'inclusion}
\[
\Fibdisc\, X \longrightarrow \Cat/X
\]
\textit{y commute, et en particulier transforme carrés cartésiens de} $\Fibdisc\, X$ \textit{en carrés cartésiens de} $\Cat/X$\textit{, ou ce qui revient au même, de} $\Cat$\textit{. De plus, en vertu de l'équivalence}
\[
\varphi_X : X^{\wedge} \xrightarrow{\ \sim\ } \Fibdisc\, X \hookrightarrow \Cat/X,
\]
\textit{le foncteur $\varphi_X$ de $X^{\wedge}$ dans} $\Cat/X$ \textit{est exacte à gauche, en particulier il transforme carrés cartésiens de $X^{\wedge}$ en carrés cartésiens de} $\Cat/X$\textit{, ou encore de} $\Cat$\textit{.} % [sic] exacte -> exact, agreement with "le foncteur"

\msp{page 165}
Notons aussi qu'étant exact à gauche, le foncteur $\varphi_X$ de $X^{\wedge}$ dans Cat$/X$ transforme monomorphismes en monomorphismes. D'autre part :

\textbf{Lemme 2.} \textit{Les monomorphismes qui sont des isomorphismes locaux sont exactement les immersions ouvertes. Donc dans} $\Fibdisc\, X$\textit{, les monomorphismes sont les immersions ouvertes, et le foncteur $\varphi_X : F \longmapsto X/F$ transforme monomorphismes en immersions ouvertes.}

Revenons alors au diagramme cocartésien $(*)$ (p. 161), et au diagramme correspondant dans Cat (et même dans $\Fibdisc\, X$)
\[
\begin{tikzcd}
X/G \ar[r, hook, "X/i"] \ar[d, "X/p"'] & X/F \ar[d, "X/q"] & T \ar[l, hook', dotted] \ar[d, dotted] \\
X/G' \ar[r, hook, "X/i'"] & X/F' & T'. \ar[l, hook', dotted]
\end{tikzcd}
\]
Comme $(*)$ est cartésien, le carré précédent est cartésien, par ce qui précède. Je veux prouver qu'il est également cocartésien
\msp{page 166}
(dans Cat$/X$, ou dans Cat, cela revient au même), et pour ceci j'applique le corollaire 2 de la proposition 15 (p. 161), compte tenu que $X/q$ est un isomorphisme local, et $X/i'$ (étant un isomorphisme local \texttt{[et un]} monomorphisme) une immersion ouverte. Il suffit de prouver que, $T'$ et $T$ étant les sous-catégories fermées complémentaires des ouverts $X/G'$ de $X/F'$ et de $X/G$ de $X/F$, que $T \longrightarrow T'$ est \texttt{[un]} isomorphisme.

\texttt{[Maintenant, $F$ et $G$, et $F'$ et $G'$, mais pas $p$ et $q$, échangent leur rôles.]} % [sic] leur rôles -> leurs rôles, agreement

La vérification est immédiate : $T$ est la catégorie des flèches $x \longrightarrow G$ dans $X^{\wedge}$ (de source $x$ dans $X$) qui ne se factorisent pas par $G$ \texttt{[plutôt par $F$]}, $T'$ celle des flèches $x \longrightarrow G'$
% foot: Version 7 décembre 2002

\end{document}
